% problem-000020 9 氨基乙醇酰化选择性
\section*{9. 氨基乙醇酰化选择性}
\textit{下列为分问。}

\subsection*{9-1}
氨基⼄醇在盐酸中与⼄酸酐反应如下：

（图：）

当此反应在 $\mathrm { K _ { 2 } C O _ { 3 } }$ 中进⾏，得到了另⼀个⾮环状的产物B；化合物A在 $\mathrm { K _ { 2 } C O _ { 3 } }$ 作⽤下也转化为化合物 $\mathbf { B } _ { \mathrm { c } }$ 。9-1-1 画出化合物B的结构简式；

\subsection*{9-1}
-2 为什么在HCl作⽤下，氨基⼄醇与⼄酸酐反应⽣成化合物A；⽽在 $\mathrm { K _ { 2 } C O _ { 3 } }$ 作⽤下，却⽣成化合物B？9-1-3 为什么化合物A在 $\mathrm { K _ { 2 } C O _ { 3 } }$ 作⽤下转化为化合物B？

\subsection*{9-2}
某同学设计如下反应条件，欲制备化合物C。但反应后实际得到其同分异构体E。

（图：）

\subsection*{9-2}
-1 画出重要反应中间体D及产物E的结构简式。

在下⾯的反应中，化合物F与⼆氯亚砜在吡啶-⼄醚溶液中发⽣反应时未能得到氯代产物，⽽是得到了两种含有碳碳双键的同分异构体G和H；没有得到另⼀个同分异构体J。

（图：）

\subsection*{9-2}
-2 画出G、H以及J的结构简式。

\subsection*{9-2}
-3 解释上述反应中得不到产物J的原因。
